\documentclass[12pt]{amsart}


\usepackage{amssymb}
\usepackage{amsthm}
\usepackage{amsmath}
\usepackage{amsxtra}
\usepackage{mathrsfs}
\usepackage{enumerate}
\usepackage{colonequals}

\setlength{\hfuzz}{4pt}

\newenvironment{problab}[1]
{\noindent\textbf{Problem #1}.}
{\vskip 6pt}

\newenvironment{exolab}[1]
{\noindent\textbf{Exercise #1}.}
{\vskip 6pt}


\theoremstyle{remark}
\newtheorem*{soln}{Solution}

\newcommand{\C}{\mathbb C}
\newcommand{\D}{\mathfrak D}
\renewcommand{\H}{\mathfrak H}
\newcommand{\F}{\mathbb F}
\renewcommand{\P}{\mathbb P}
\newcommand{\PP}{\mathbb P}
\newcommand{\Q}{\mathbb Q}
\newcommand{\R}{\mathbb R}
\renewcommand{\SS}{\mathbb S}
\newcommand{\Z}{\mathbb Z}

\renewcommand{\Im}{\operatorname{Im}}
\renewcommand{\Re}{\operatorname{Re}}

\newcommand{\eps}{\epsilon}

\usepackage{fullpage}

\DeclareMathOperator{\Aut}{Aut}
\DeclareMathOperator{\Hom}{Hom}
\DeclareMathOperator{\ord}{ord}

\usepackage{mathpazo}

\begin{document}

\title{Belyi Maps and Dessins d'Enfants \\ Homework \#4}
\date{\today}

%\thanks{\emph{Date:} Monday, 11 February 2013}

\maketitle

\begin{exolab}{4.1}
  Use the fact that $\chi(\SS^2) = 2$ to show that there are only five regular polyhedra.  [\textit{Hint}: Consider subdivisions of the sphere into $n$-gons such that exactly $m$ edges meet at each vertex, $m ,n \geq 3$.  Show that $nf = 2e = mv$, so that
  $$
    \frac{1}{e} + \frac{1}{2} = \frac{1}{m} + \frac{1}{n}
  $$
  and consider the solutions with $n = 3, 4, 5$ and $n \geq 6$ in turn.]
\end{exolab}


\begin{exolab}{4.2}
  Let $E: Y^2 Z = F(X,Z)$ be an elliptic curve given in Weierstrass form. Let $x = X/Z$ and $y = Y/Z$ be affine coordinates on the affine open $U$ where $Z \neq 0$ and let $\omega = \frac{dx}{y}$. We showed in class that $\omega$ is holomorphic on $U$. Let $\infty = [0:1:0]$.
  \begin{enumerate}[(a)]
    \item 
      Show that $\ord_\infty(x) = -2$ and $\ord_\infty(y) = -3$. [\textit{Hint}: What are the degrees of $x$ and $y$? How many points in their respective fibers above $\infty = [1:0] \in \PP^1$?]
    \item
      Show that $\omega$ is also holomorphic at $[0:1:0]$ and conclude that it is a holomorphic differential on all of $E$. [\textit{Hint}: Compute $\ord_\infty(\omega)$ using the fact that $\ord_\infty(dx) = \ord_\infty(x) - 1$.]
  \end{enumerate}
\end{exolab}

%\begin{exolab}{4.3}
%  % HW? y^2 = x^5 - 1; show dx/y and x*dx/y are holomorphic
%  Let $C: Y^2 = X^5 Z - Z^6$ be a hyperelliptic curve, given as a subset of $\P(1,3,1)$. Let $x = X/Z$ and $y = Y/Z^3$ be affine coordinates on the affine open subset $U_2$ where $Z \neq 0$. Show that the differentials
%  \begin{equation*}
%    \omega_1 \colonequals \frac{dx}{y}, \qquad \omega_2 \colonequals x \frac{dx}{y}
%  \end{equation*}
%  are holomorphic on all of $C$.
%\end{exolab}

\end{document}
